% NAF 15383, batch 1 — Gallica views 1–20.
% Pass: Opus 5 (claude-opus-5), 2026-09-16 — first pass, unchecked against the
% views by a human.
%
% What the twenty views hold. Five carry writing; the other fifteen are the
% binding, its case, and blank leaves.
%
%   v. 1–8    Binding and case. Front board; marbled contreplat; the marbled
%             chemise photographed open (v. 3) and unfolded flat (v. 4); the
%             inner board with the printed library label « FR / NOUV. ACQ. /
%             15383 / 1516·IV » (v. 4 and v. 5); then three blank guard leaves
%             (v. 6–8), examined at full contrast and carrying nothing.
%             Not transcribed.
%   v. 9–10   An oblong leaf, separate from the letter and not mentioned in the
%             catalogue entry: a quittance for the rentes of the Hôtel de Ville
%             de Paris, passed at Clermont on the twenty-seventh of June 1674
%             by « Dame Gilberte Pascal veufve et heritiere testamentaire de
%             deffunt Mr Me Florin Perier », and signed by her — « G. Pascal ».
%             Its verso (v. 10) carries a certificate of 4 July 1674 by Antoine
%             Ferret, bourgeois de Paris, with his paraph. The IIIF manifest
%             labels views 9–16 « NP »; these two are a second piece, in
%             landscape format, thirty-one years later than the letter.
%   v. 11–12  Blank leaves.
%   v. 13     The letter, first page.
%   v. 14     Its verso: blank, the recto showing through, with the round stamp
%             « BIBLIOTHÈQUE NATIONALE MSS ». No placeholder is written for it.
%   v. 15     The letter, second written page, with the signature; and, in the
%             blank space left in the middle of the page, a note in a second
%             hand.
%   v. 16     The address panel.
%   v. 17–18  Blank leaves.
%   v. 19–20  The chemise unfolded and the back board. Not transcribed.
%
% The letter, and who writes to whom. It begins at the head of v. 13 — « De
% Rouen Ce Samedy dernier Januier 1643 », then « Ma Chere Seur » — runs to the
% foot of that page in mid-sentence (« Tellement que Nous Sommes »), and is
% taken up at the head of v. 15 (« En Vne peine dont Je te prie de Nous tirer
% au plustost »). It ends on v. 15 with « Ma Chere Seur » set large at the left
% and, at the lower right, « Vostre tres humble et tres affectionné serviteur
% et frere » over the signature « Pascal ». The writer names himself the
% brother, speaks of « mon Pere » in the third person, and writes from Rouen;
% the address on v. 16 is to « Mademoiselle Perier La Conseillere a Clermont ».
% Both the date and both signatures are legible.
%
% The second hand. Into the space left blank between the end of the letter and
% its closing formula, a smaller, much faster hand has written a note of
% thirteen lines beginning « Ma bonne fille Mignonne », complaining of having
% no leisure, and closing « Vre bon Pere et tres affectionné amy » over a
% second signature « Pascal ». It is given here after the brother's text, in
% its own paragraph, with a \note{} saying where it stands on the leaf. This
% hand is markedly harder to read than the letter's: several of its words are
% left \ill{}.
%
% The hand of the letter. A clear, wide, well-spaced cursive with generous
% flourishes. The long s is transcribed as s throughout, said once here. The
% final r is often a small closed loop that reads like an o (« Seur »,
% « Pere »). Capitals are used freely inside the sentence (Ce, Cela, Croy,
% Nous, Une, Sepmaines): kept as written. « que » is written as a q with a
% looped flourish and is set « que ». Abbreviations are kept: « lres » for
% « lettres » (the res surmounted by a bar), « Mr », « Mrs », « Mde », « Vre »,
% « ordre » for « ordinaire », « Nores » for « notaires », « apnt » for « a
% present ». The spelling of 1643 stands unflagged — escripre, escripvois,
% asseuré, sepmaines, doubté, estonnois, plustost, tems, scavoir, Januier,
% Moy for « mois », veufve, moys, aud, lesd, lad.
%
% Numbers on the leaves. No foliation of the library is legible anywhere in the
% batch — not on v. 13, not on v. 15, not on v. 9 — so no \folio{} is written.
% What the leaves do carry: the modern pencil shelfmark « N. a. fr. 15383 »,
% underlined, at the top left of v. 13 and at the foot of v. 9; the oval
% acquisition stamp « ACHAT No 24199 » on both; a pencilled acquisition or
% dealer's note at the foot of v. 9 (a name, then « 34. 226tt 10/ »); and, in
% the left margin of v. 9, three clerk's marks in ink, among them the figure
% 84. Inside the two documents the figures are the writers' own: the year 1643
% on v. 13, and on v. 9–10 the sums and dates of the rente (133 livres 4 sols,
% 888 livres, 1635, 1665, 1674) written out in words, with one sum on v. 10 in
% the old score notation.
%
% The catalogue gives no dating for this volume, so \dating{} is left empty.

\documentclass[11pt,a4paper]{article}
% The preamble shared by every transcription of the mathematicians' manuscripts in Gallica.
%
% It rests on one idea: **uncertainty compounds, it does not resolve.** A
% transcription that smooths over a doubtful word has destroyed the only thing
% it was made for — a reader must be able to tell, at every point, what was on
% the page and what was supplied. The apparatus macros below are therefore the
% critical apparatus, not decoration.
%
% The same commands are recognised by `scripts/render.mjs`, which produces the
% site's reading view, and by `scripts/tei.mjs`. A macro added here must be
% added there too, or rendering fails loudly — which is the wanted behaviour.
%
% Comments here are English, like the rest of the repository. The documents
% this preamble sets are French, and that is deliberate: see the skills.

\ProvidesPackage{germain}[2026/09/15 Transcriptions of mathematicians' manuscripts in Gallica]

\usepackage{amsmath,amssymb,amsthm}
\usepackage{mathrsfs}
% Kept from the parent preamble so the renderer's diagram support stays
% available; these manuscripts draw no commutative diagrams, and nothing here uses it.
\usepackage{tikz-cd}
\usepackage[english,french]{babel}
\usepackage{fontspec}

% --- Greek ------------------------------------------------------------------
% The text face stays fontspec's default, Latin Modern, which has no Greek:
% XeTeX drops every glyph it lacks with a line in the log and nothing on the
% page, and the Greek words the manuscripts quote vanished from the PDFs. So
% the Greek and Greek Extended blocks get a XeTeX character class of their own,
% and entering or leaving a run of it switches the family to CMU Serif — the
% same Computer Modern design, with polytonic Greek — and back. Only the
% family changes: a Greek word in \textit or \textbf keeps its shape and series.
% The transitions to and from every other class carry the tokens that class
% has towards an ordinary letter, so French spacing before « ; » or inside
% « » still holds after a Greek word. No grouping, so no brace or \endgroup
% can come out unbalanced; maths is left alone, since XeTeX inserts nothing
% there. Kept identical in `exercices.sty`.
\makeatletter
\newfontfamily\germainGreekfont{cmun}[
  Extension=.otf, NFSSFamily=germaingreek,
  UprightFont=*rm, ItalicFont=*ti, SlantedFont=*sl,
  BoldFont=*bx, BoldItalicFont=*bi, BoldSlantedFont=*bl]
\newXeTeXintercharclass\germain@greek
\def\germain@greekrange#1#2{%
  \@tempcnta=#1\relax
  \loop
    \XeTeXcharclass\@tempcnta=\germain@greek
    \ifnum\@tempcnta<#2\advance\@tempcnta\@ne\repeat}
\germain@greekrange{"0370}{"03FF}
\germain@greekrange{"1F00}{"1FFF}
\def\germain@greekprev{\rmdefault}
\def\germain@greekin{%
  \edef\germain@greekprev{\f@family}\fontfamily{germaingreek}\selectfont}
\def\germain@greekout{\fontfamily{\germain@greekprev}\selectfont}
\def\germain@greekjoin{%
  \XeTeXinterchartokenstate=\@ne
  \@tempcnta=\z@
  \loop
    \ifnum\@tempcnta=\germain@greek\else
      \XeTeXinterchartoks\germain@greek\@tempcnta=\expandafter{%
        \expandafter\germain@greekout\the\XeTeXinterchartoks\z@\@tempcnta}%
      \XeTeXinterchartoks\@tempcnta\germain@greek=\expandafter{%
        \the\XeTeXinterchartoks\@tempcnta\z@\germain@greekin}%
    \fi
    \ifnum\@tempcnta<4095\advance\@tempcnta\@ne\repeat}
% babel-french sets its own classes, and saves and restores the interchar
% state, each time a language is selected: join again after it does.
\addto\extrasfrench{\germain@greekjoin}
\addto\extrasenglish{\germain@greekjoin}
\AtBeginDocument{\germain@greekjoin}
\makeatother

\usepackage{geometry}
\geometry{a4paper,margin=25mm,marginparwidth=28mm}
\usepackage{marginnote}
\usepackage{xcolor}
\usepackage[normalem]{ulem}
\setlength{\emergencystretch}{3em}
\usepackage{hyperref}
\hypersetup{colorlinks=true,linkcolor=black,urlcolor=black,pdfborder={0 0 0}}

% --- Batch metadata ---------------------------------------------------------
% Taken as they stand by the rendering script, which shows them at the head of
% the reading view. They fix what the file claims to cover. The macro names are
% the parent project's, so that the scripts stay shared; \folder here means the
% volume's slug — `fr-9115` — and \pages counts Gallica views.

\newcommand{\folder}[1]{\gdef\grothFolder{#1}}
\newcommand{\batch}[1]{\gdef\grothBatch{#1}}
\newcommand{\pages}[2]{\gdef\grothFirst{#1}\gdef\grothLast{#2}}
\newcommand{\foldertitle}[1]{\gdef\grothFoldertitle{#1}}

% The holder's own shelfmark, printed as they write it: « Français 9115 ».
\newcommand{\shelfmark}[1]{\gdef\grothShelfmark{#1}}

% The dating, copied verbatim from the holder's catalogue — never inferred
% here. « XIXe siècle » is the BnF's; a pass that dates a leaf from its content
% says so in a \note{}, as its own inference.
\newcommand{\dating}[1]{\gdef\grothDating{#1}}

% The status notice, printed in the title block and repeated as a diagonal
% watermark on every page of the PDF. The manuscripts are in the public
% domain; what this line declares is not a right but a quality — an
% unverified first pass by a machine. The skills write the line; a file
% without it gets no watermark, so leaving it out is visible in the output.
\newcommand{\watermark}[1]{\gdef\grothWatermark{#1}}

\gdef\grothFolder{?}\gdef\grothBatch{}\gdef\grothFirst{?}\gdef\grothLast{?}%
\gdef\grothFoldertitle{}\gdef\grothShelfmark{}\gdef\grothDating{}\gdef\grothWatermark{}

% --- The résumé -------------------------------------------------------------
\newenvironment{resume}
  {\small\setlength{\parskip}{0.45em}}
  {\par\medskip\hrule\medskip}

% --- Keywords ---------------------------------------------------------------
% In English, closing the modernised reading's résumé: the modern vocabulary
% under which to search for what the leaves construct. The manifest extracts
% them to tag the volume — this line is the single source of the tags.
\newcommand{\keywords}[1]{\par\smallskip\noindent
  {\footnotesize\textsc{Keywords} --- \textit{#1}}\par}

% --- The critical apparatus -------------------------------------------------

% The Gallica view number, in the margin. This is the anchor that lets a
% disagreement over a formula be settled by looking at *one* image rather than
% a volume of 750. The site uses it to turn the facsimile as the transcription
% is scrolled. A view is one scanned image: usually one side of a leaf.
\newcommand{\page}[1]{%
  \marginnote{\footnotesize\color{gray!70}\textsf{v.\,#1}}%
  \label{page:#1}\ignorespaces}

% The BnF's pencilled foliation, where the leaf carries it and a pass can read
% it: « 198r ». This is what the literature cites — « ff. 198r–208v » — and
% the one line that turns a citation into a view. Recorded, never inferred:
% a folio a pass cannot read is not written.
\newcommand{\folio}[1]{%
  \marginnote{\footnotesize\color{gray!50}\textsf{f.\,#1}}\ignorespaces}

% The modernised reading's coarser counterpart to \page{}: which views a
% section is drawn from.
\newcommand{\pagerange}[2]{%
  \marginnote{\footnotesize\color{gray!70}\textsf{v.\,#1--#2}}%
  \ignorespaces}

% Illegible: never guessed.
\newcommand{\ill}{\textcolor{red!60!black}{[\,\ldots\,]}}

% A doubtful reading: the word is given, and flagged as uncertain.
\newcommand{\uncertain}[1]{\uline{#1}}

% An editorial addition — a missing letter, an implied word.
\newcommand{\add}[1]{\textcolor{blue!50!black}{[#1]}}

% Struck out by the writer. What was crossed out is part of the document.
\newcommand{\struck}[1]{\sout{#1}}

% The transcriber's note, on the state of the page or the mathematics.
\newcommand{\note}[1]{\par\smallskip{\small\color{gray!60!black}\textit{[#1]}}\par\smallskip}

% A marginal note of hers, distinct from the body of the text.
\newcommand{\marginal}[1]{\par\smallskip\noindent\hspace{1em}%
  {\small\color{gray!60!black}#1}\par\smallskip}

% --- Title ------------------------------------------------------------------
% The title block is French, like the documents themselves.

\AddToHook{shipout/background}{%
  \ifx\grothWatermark\empty\else
    \begin{tikzpicture}[remember picture, overlay]
      \node[rotate=52, scale=3.4, align=center, text=gray!18,
            font=\sffamily\bfseries]
        at (current page.center) {\grothWatermark};
    \end{tikzpicture}%
  \fi}

\AtBeginDocument{%
  \begin{center}
    {\large\bfseries Manuscrits de mathématiciens (Gallica) ---
      \ifx\grothShelfmark\empty\grothFolder\else\grothShelfmark\fi}\\[2pt]
    {\itshape\grothFoldertitle}\\[4pt]
    {\small vues \grothFirst--\grothLast
      \ifx\grothBatch\empty\else\ (lot \grothBatch)\fi}\\[2pt]
    \ifx\grothDating\empty\else
      {\small\color{gray!60!black}Datation du catalogue : \grothDating}\\[2pt]
    \fi
    \ifx\grothWatermark\empty\else
      {\small\bfseries\color{red!45!gray}\grothWatermark}\\[2pt]
    \fi
    {\footnotesize\color{gray!60!black}Transcription établie d'après les images IIIF de Gallica.
     Source gallica.bnf.fr / Bibliothèque nationale de France.\\
     Les lectures incertaines sont soulignées, les illisibles notées [\,\ldots\,],
     les ajouts entre crochets ; \textsf{v.} numérote les vues de Gallica,
     \textsf{f.} la foliotation au crayon de la BnF.}
  \end{center}
  \vspace{1em}}

\endinput


\folder{naf-15383}
\shelfmark{NAF 15383}
\batch{1}
\pages{1}{20}
\dating{}
\watermark{Lecture automatique\\non vérifiée}
\foldertitle{Lettre de Blaise Pascal à sa sœur Gilberte, 31 janvier 1643.XVIIe s.}

\begin{document}

\page{9}
\note{Feuillet à l'italienne, indépendant de la lettre : une quittance de rente sur l'Hôtel de Ville de Paris, passée à Clermont le 27 juin 1674. Le titre et le cartouche ovale sont gravés ; le corps est d'une main de notaire. Dans la marge de gauche, trois marques à l'encre — un f de scribe, le chiffre 84 suivi d'un point, et un signe en forme de 8 — puis le cachet ovale « ACHAT No 24199 ». Au bas du feuillet, au crayon et d'une main moderne, la cote « N. a. fr. 15383 » soulignée, et une note d'acquisition dont le nom n'a pu être lu. Aucune foliation de la bibliothèque.}

Quittance pour les Rentes de l'Hostel de Ville de Paris.
\note{titre gravé, coupé en son milieu par un cartouche ovale portant, en cercle, « GENERALITE DE PARIS », et une fleur de lys.}

Deux sols.

1674. derniers six moys.
\note{en marge de gauche, en haut, souligné.}

Par devant les Notaires Royaux establis en la Ville de Clermont fut presente Dame Gilberte Pascal veufve et heritiere testamentaire de deffunt M\textsuperscript{r} M\textsuperscript{e} Florin Perier vivant Cons\textsuperscript{er} du Roy en sa Cour des Aydes de Clermont ferrand, Laquelle de son bon gré a Confessé avoir receu de
\note{la ligne s'arrête là : le nom du payeur est resté en blanc.}

La somme de Cent trente trois livres quatre sols ordonnée estre laissée en fonds dans l'Estat du Roy de la Ferme generale des Aydes et Entrées pour les six derniers moys de l'année presente mille six cent soixante \& quatorze A cause de huit cent quatre vingt huit livres de rente constituée sur les tailles par contrat du second Janvier mille six cent trente cinq, Et a present assignée sur lesd Aydes et Entrées suivant la Declaration de sa Majesté du treiziesme Janvier mille six cent soixante cinq. De laquelle somme de Cent trente trois livres quatre sols lad Dame quitte lesd Sieur payeur et tous autres. Ainsy \&c. A peine \&c. Oblige \&c. Renonce \&c. Soumet \&c. Fait et passé aud Clermont Estude de l'un desd Notaires le vingt septiesme jour de Juin. mille six cent soixante quatorze apres midy. Lad Dame Confessante a signé avec Lesd No\textsuperscript{res} J'approuve Le mot de six \ill{}
\note{« six derniers moys » : « six » est ajouté au-dessus de la ligne, ce que la mention finale du notaire approuve. Après « Le mot de six », la fin de l'approbation, en petite cursive rapide, n'a pu être lue. Un long trait de remplissage suit « vingt septiesme ».}

Quittance de Cent trente trois livres quatre sols.
\note{ligne de dos, tracée au bas du texte et en partie recouverte par les paraphes.}

Camus
\note{signature, suivie de « no\textsuperscript{re} royal a Clermont » et d'un grand paraphe.}

G. Pascal.
\note{signature autographe de Gilberte Pascal, au milieu de la ligne des signatures ; le cachet rond « BIBLIOTHÈQUE NATIONALE MSS » est apposé juste au-dessous.}

Guillotteau
\note{signature, suivie de « no\textsuperscript{re} royal » et d'un grand paraphe.}

\page{10}
\note{Verso du feuillet précédent ; le recto y transparaît en entier, à l'envers. Le certificat est écrit en haut à droite.}

Je Soubz\textsuperscript{ne} Antoine \uncertain{Ferret} bourgeois de paris, Certiffie a tous quil appartiendra que la quittance de l'autre part est Veritable a Paris ce quat\textsuperscript{e} Juillet \uncertain{mil vj}\textsuperscript{c} Soix\textsuperscript{te} quatorze
\note{le millésime est écrit en abrégé, deux ou trois lettres surmontées d'un c : la lecture « mil vj\textsuperscript{c} » est offerte par le sens, « Soix\textsuperscript{te} quatorze » étant net. Le nom est donné d'après la signature qui suit ; l'initiale pourrait aussi bien se lire L.}

\uncertain{Ferret}
\note{signature, suivie d'un très grand paraphe en boucles.}

\note{Au milieu du feuillet, à gauche de la signature, un chiffre à l'ancienne notation par vingtaines dont seule la fin se lit : \ill{}\textsuperscript{xx} xiij. Il s'accorderait avec les cent trente-trois livres de la quittance, mais la pièce ne le dit pas. En bas à droite, tête-bêche, quatre lignes très pâles — un millésime 1674, le mot « escheue », un mois et des chiffres — qui n'ont pas été lues.}

\page{13}
\note{Première page de la lettre. En haut à gauche, au crayon et d'une main moderne, la cote « N. a. fr. 15383 » soulignée, et le cachet ovale « ACHAT No 24199 ». Le cachet rond « BIBLIOTHÈQUE NATIONALE MSS » est apposé dans la dernière ligne du texte. Aucune foliation de la bibliothèque n'est lisible sur le feuillet.}

De Rouen Ce Samedy dernier Januier 1643
\note{en tête, à droite, sur deux lignes. Les deux derniers chiffres du millésime sont tracés en grand et descendent sous la ligne : le 4 est barré de sa traverse, le 3 est formé de deux boucles. Le « 16 » est net.}

Ma Chere Seur

Je ne doubte pas que vous n'ayez esté bien en peine du long temps qu'il y a que vous n'avez receu de Nouvelles de Ces quartiers Jcy. Mais Je Croy que vous vous serez bien doubté que le voyage des Esleuës en a esté la Cause. Comme en effect sans Cela Je n'aurois pas manqué de vous Escripre plus souvent. J'ay a te dire que M\textsuperscript{rs} Les Commissaires estans a Gizors mon Pere me fist aller faire Un tour a Paris ou Je trouvé Une lettre que tu m'escripvois ou tu me mandois que tu t'estonnois de Ce que Je te reproche que tu n'escripvois pas assez souvent; ou tu me dis que tu escripts a Rouen toutes les Sepmaines Une fois.
\note{le passage du « vous » au « tu » se fait au milieu d'une ligne, à « J'ay a te dire », et ne revient plus. « Gizors » est écrit d'un seul trait, le z bouclé.}

Jl est bien asseuré Si Cela est que les l\textsuperscript{res} se perdent Car Je n'en \uncertain{reçois} pas toutes les trois Sepmaines Une
\note{« l\textsuperscript{res} » : l'abréviation ordinaire de « lettres », le l bouclé et le groupe res surmonté d'un trait. Le verbe est écrit vite et la cédille n'est pas sûre.}

Estant retournez a Rouen J'y ay trouvé Une lettre de M\textsuperscript{r} Perier qui mande que tu es \uncertain{tres} malade Jl ne mande point Si ton mal est dangereux ni Si tu te portes mieux, et Jl s'est passé Un ord\textsuperscript{re} depuis sans avoir rien de lettres Tellement que Nous Sommes
\note{le mot entre « es » et « malade » est bref, écrit en partie au-dessus de la ligne et surmonté d'un trait d'abréviation, sur une pliure du papier : la lecture « tres » est offerte avec doute. « ord\textsuperscript{re} » est l'abréviation usuelle de « l'ordinaire », le courrier. La page s'arrête ici, en pleine phrase ; le bas du feuillet est blanc et taché. La phrase se poursuit à la vue 15.}

\page{15}
\note{Deuxième page écrite de la lettre : elle reprend la phrase laissée en suspens à la vue 13. Le feuillet porte deux mains. Celle de la lettre occupe les sept premières lignes, puis reparaît en bas pour le « Ma Chere Seur » et la formule finale ; entre les deux, dans le blanc laissé au milieu de la page, une seconde main, plus petite et beaucoup plus rapide, a écrit un billet de treize lignes. Ce billet est donné plus bas. Aucune foliation.}

En Vne peine dont Je te prie de Nous tirer au plustost mais Je Croy que la prie que Je te fais Jcy Sera Jnutile Car avant que tu ayes receu Cette lettre Jcy J'espere que \struck{ou tu} Nous aurons receu lettres ou de toy ou de Monsieur Perier. Le depart \ill{} du fascheux du Moy.
\note{« la prie » est suivi d'un long trait horizontal, peut-être pour « priere ». Le mot qui suit « Le depart » compte cinq ou six jambages, traverse une pliure du papier, et n'a pas été lu ; « du fascheux du Moy » est net, « Moy » étant ici le mois.}

Si Je scavois quelque Chose de Nouveau Je te le ferois scavoir. Je Suis

Ma Chere Seur
\note{écrit en grand, à gauche, au bas de la page.}

Vostre tres humble et tres affectionné serviteur et frere
\note{en bas à droite, sur deux lignes.}

Pascal
\note{signature, en grand, avec un paraphe.}

\note{Ce qui suit est le billet de la seconde main, écrit au milieu de la page, dans le blanc laissé entre la fin de la lettre et sa formule finale. L'écriture est petite, très cursive, à longs traits horizontaux, et par endroits pâlie ; plusieurs mots n'ont pu être lus.}

Ma bonne fille Mignonne Si \uncertain{Vous} \ill{} d'escrire N'y ayant aucun loisir Car Je n'ay Jamais esté dans l'embarras a la des\textsuperscript{me} \ill{} de ce que J'y suis a p\textsuperscript{nt}, Je ne scaurois L'estre d'avantage a moins d'y avoir trop. Jl y a \struck{trois} quatre Moy que Je ne suis pas couché six fois devant deux heures apres minuict.
\note{la fin de la première ligne, quatre ou cinq mots entre « Si Vous » et « d'escrire », n'a pas été lue. « des\textsuperscript{me} » est suivi d'un mot d'une dizaine de jambages, non lu. « quatre » est écrit au-dessus de « trois » biffé. « Moy » est ici le mois.}

Je Vous envoy \ill{} Une l\textsuperscript{re} de \ill{} sur Ce subject de la l\textsuperscript{re} derniere touchant les \uncertain{bruits} du Mariage de Monsieur \uncertain{Despeaux} Mais Je n'ay Jamais sceu \uncertain{que Ce fust} de L'archevesque.
\note{entre « envoy » et « Une l\textsuperscript{re} », deux ou trois mots non lus. Le nom qui suit « de » compte huit ou neuf lettres, commence par une majuscule qui se lit B ou R et paraît finir en -bre : il n'est pas donné. « Despeaux » est offert avec doute, le milieu du mot étant surchargé.}

Pour Nouvelles La fille de Monsieur de Paris M\textsuperscript{de} \uncertain{du Comptes} Mariée a Monsieur de \uncertain{Montiffer} aussy M\textsuperscript{r} \uncertain{du Comptes} est decedée Comme aussy la fille de \uncertain{Belair} mariée au petit \uncertain{Lambert} Je suis tousjours
\note{les noms propres de ce paragraphe sont écrits très vite et sans reprise ; ils sont donnés tels qu'ils se présentent, et marqués douteux.}

V\textsuperscript{re} bon Pere et tres affectionné amy
\note{formule finale du billet, à droite.}

Pascal
\note{seconde signature, au milieu de la page, en grand ; le cachet rond de la bibliothèque la chevauche.}

V\textsuperscript{re} petit a couché \uncertain{ceste} nuict Il se porte Dieu grace tres bien
\note{trois lignes serrées dans la marge inférieure gauche, de la même seconde main, en regard de la formule finale. Un mot bref à la fin de la première ligne n'a pas été lu.}

\page{16}
\note{Verso du second feuillet : le panneau d'adresse, écrit de la main de la lettre au milieu de la page, le feuillet ayant été plié. Le texte de la vue 15 transparaît. Aucune foliation, aucune marque postale lisible.}

A Mademoiselle

Mademoiselle Perier La

Conseillere

a Clermont
\note{sur quatre lignes, la seconde et la troisième en retrait. « La Conseillere » désigne la femme du conseiller.}

\end{document}
